\subsection{Configuração Experimental}

O desempenho dos modelos é avaliado por meio do F1-Score médio por categoria, reportado como média e desvio padrão em 5 \textit{folds} com divisão estratificada, utilizando 80\% dos dados para treino e 20\% para validação, preservando a proporção entre categorias.

\begin{table}[htbp]
\centering
\caption{Total de \textit{check-ins}, POIs, e usuários de Flórida, Califórnia e Texas.}
\label{tab:dataset}
\begin{tabular}{lccc}
\hline
Estados & \textit{Check-ins} & Pontos de Interesse & Usuários \\
\hline
Florida & 990.518 & 65.009 & 20.301 \\
California & 2.535.573 & 148.314 & 36.106 \\
Texas & 3.355.419 & 135.570 & 37.522 \\
\hline
\end{tabular}
\end{table}

Os experimentos foram conduzidos com o \textit{dataset} Gowalla \cite{10.1145/2661829.2662002} nos estados da Flórida, Califórnia e Texas, cujas estatísticas são apresentadas na Tabela \ref{tab:dataset}. A escolha desses estados se deve ao elevado volume de \textit{check-ins} e à presença de padrões urbanos e distribuições espaciais distintas, o que permite avaliar o comportamento dos modelos em cenários com diferentes níveis de densidade e dispersão geográfica.

\begin{figure}[htbp]
    \centering
    \includegraphics[width=\textwidth]{imagens/subáreas/distribuicao_estados.png}
    \caption{Distribuição espacial de POIs das categorias Food (vermelho) e Shopping (laranja) na sub-região mais densa de Flórida, Califórnia e Texas. As sub-regiões foram selecionadas por densidade em \textit{grid}, com cerca de 100 POIs por região.}
    \label{fig:distribuicao_espacial}
\end{figure}

\begin{table}[ht]
\centering
\caption{\textnormal{F1-Score Médio (\%) por modelo e estado para a tarefa de Classificação de Categoria de POI.}}
\label{table:poi_category}
\vspace{0.3cm}
\resizebox{\textwidth}{!}{%
\begin{tabular}{cl|c|c|c|}
\hline
\multicolumn{2}{c|}{} & \textbf{MTLNet} & \textbf{ST-MTLNet$_{SIREN}$} & \textbf{ST-MTLNet$_{Sphere2Vec-M}$} \\ \hline
\multirow{7}{*}{\textbf{Florida}} & Community & 51,86 ± 0,73 & \textbf{70,00 ± 0,81} & 69,84 ± 0,98 \\
 & Entertainment & 41,24 ± 1,83 & \textbf{64,45 ± 1,82} & 63,92 ± 1,61 \\
 & Food & 55,47 ± 1,55 & 72,92 ± 0,59 & \textbf{73,22 ± 0,66} \\
 & Nightlife & 32,59 ± 1,96 & \textbf{62,60 ± 1,96} & 62,44 ± 2,00 \\
 & Outdoors & 47,71 ± 1,60 & 65,64 ± 1,73 & \textbf{66,35 ± 1,51} \\
 & Shopping & 62,96 ± 0,62 & \textbf{77,48 ± 0,36} & 77,47 ± 0,72 \\
 & Travel & 45,49 ± 1,20 & \textbf{64,89 ± 1,20} & 64,77 ± 1,54 \\
\hline
\multirow{7}{*}{\textbf{California}} & Community & 53,22 ± 0,37 & \textbf{70,21 ± 0,17} & 70,08 ± 0,44 \\
 & Entertainment & 40,89 ± 1,32 & \textbf{63,60 ± 0,83} & 63,21 ± 0,76 \\
 & Food & 60,68 ± 0,69 & \textbf{75,78 ± 0,39} & 75,74 ± 0,19 \\
 & Nightlife & 26,81 ± 1,71 & 60,59 ± 1,35 & \textbf{60,78 ± 1,63} \\
 & Outdoors & 51,59 ± 1,15 & \textbf{67,94 ± 1,52} & 67,37 ± 1,95 \\
 & Shopping & 60,43 ± 1,01 & 76,84 ± 0,29 & \textbf{77,00 ± 0,33} \\
 & Travel & 38,88 ± 1,32 & \textbf{63,59 ± 1,45} & 63,25 ± 0,99 \\
\hline
\multirow{7}{*}{\textbf{Texas}} & Community & 57,37 ± 0,69 & 73,45 ± 0,34 & \textbf{76,20 ± 0,46} \\
 & Entertainment & 46,69 ± 0,46 & 64,09 ± 1,24 & \textbf{67,73 ± 1,43} \\
 & Food & 54,31 ± 0,42 & 69,80 ± 0,80 & \textbf{73,10 ± 0,17} \\
 & Nightlife & 34,44 ± 1,21 & 60,62 ± 1,32 & \textbf{64,57 ± 1,03} \\
 & Outdoors & 45,99 ± 1,70 & 64,21 ± 0,42 & \textbf{68,72 ± 0,70} \\
 & Shopping & 62,70 ± 0,29 & 76,40 ± 0,44 & \textbf{79,67 ± 0,14} \\
 & Travel & 39,37 ± 1,29 & 57,53 ± 1,08 & \textbf{64,73 ± 0,72} \\
\hline
\end{tabular}}
\end{table}

\subsection{Classificação de Categoria de POI}

Os resultados da tarefa de classificação de categoria de POI são apresentados na Tabela \ref{table:poi_category}. Nos três estados avaliados, as duas variantes do ST-MTLNet superam o MTLNet original em todas as 21 combinações de categoria e estado, evidenciando que a substituição do \textit{embedding} monolítico do DGI por representações desacopladas enriquece de forma consistente a caracterização dos POIs. Os ganhos são particularmente expressivos em categorias como \textit{Nightlife} e \textit{Travel}. Em \textit{Nightlife}, por exemplo, o ST-MTLNet com Sphere2Vec-M atinge 62,44 $\pm$ 2,00 na Flórida, 60,78 $\pm$ 1,63 na Califórnia e 64,57 $\pm$ 1,03 no Texas, sempre muito acima do \textit{baseline}. Em \textit{Travel}, os melhores resultados chegam a 64,89 $\pm$ 1,20 na Flórida com SIREN, 63,59 $\pm$ 1,45 na Califórnia com SIREN e 64,73 $\pm$ 0,72 no Texas com Sphere2Vec-M.

As melhorias observadas em \textit{Food} e \textit{Shopping} reforçam que a localização geográfica carrega sinais discriminativos relevantes para a categoria semântica do POI. A Figura \ref{fig:distribuicao_espacial} mostra que essas categorias tendem a formar aglomerados densos nas sub-regiões analisadas, favorecendo a captura de padrões espaciais pelos \textit{encoders} contínuos. Embora as duas variantes propostas apresentem desempenho próximo, observa-se um comportamento complementar: o SIREN obtém mais resultados de destaque na Flórida e na Califórnia, enquanto o Sphere2Vec-M se destaca com maior frequência no Texas. Em conjunto, esses resultados mostram que a abordagem modular supera de forma consistente o paradigma monolítico do DGI nessa tarefa.

\begin{table}[ht]
\centering
\caption{\textnormal{F1-Score Médio (\%) por modelo e estado para a tarefa de Predição de Próximo POI.}}
\label{table:next}
\vspace{0.3cm}
\resizebox{\textwidth}{!}{%
\begin{tabular}{cl|c|c|c|}
\hline
\multicolumn{2}{c|}{} & \textbf{MTLNet} & \textbf{ST-MTLNet$_{SIREN}$} & \textbf{ST-MTLNet$_{Sphere2Vec-M}$} \\ \hline
\multirow{7}{*}{\textbf{Florida}} & Community & 34,33 ± 1,33 & \textbf{38,31 ± 0,90} & 37,98 ± 1,22 \\
 & Entertainment & 25,34 ± 1,07 & \textbf{31,38 ± 1,72} & 31,24 ± 2,41 \\
 & Food & 27,12 ± 2,79 & \textbf{41,91 ± 1,93} & 41,65 ± 3,15 \\
 & Nightlife & 21,33 ± 1,61 & 23,32 ± 2,74 & \textbf{23,98 ± 1,94} \\
 & Outdoors & \textbf{21,61 ± 0,99} & 21,29 ± 1,59 & 21,59 ± 1,91 \\
 & Shopping & 39,62 ± 3,40 & 44,00 ± 2,04 & \textbf{44,66 ± 2,38} \\
 & Travel & \textbf{64,47 ± 1,02} & 45,00 ± 1,10 & 44,93 ± 1,11 \\
\hline
\multirow{7}{*}{\textbf{California}} & Community & 32,14 ± 0,89 & \textbf{37,86 ± 0,20} & 37,80 ± 1,11 \\
 & Entertainment & \textbf{19,38 ± 0,94} & 17,28 ± 1,73 & 17,01 ± 1,39 \\
 & Food & 29,26 ± 1,67 & 50,42 ± 1,15 & \textbf{51,28 ± 0,57} \\
 & Nightlife & \textbf{18,24 ± 0,53} & 16,34 ± 1,95 & 15,10 ± 2,70 \\
 & Outdoors & \textbf{25,01 ± 0,81} & 24,84 ± 1,70 & 24,07 ± 2,46 \\
 & Shopping & 38,41 ± 2,64 & \textbf{39,04 ± 1,72} & 37,82 ± 1,31 \\
 & Travel & \textbf{46,05 ± 0,84} & 36,94 ± 1,70 & 37,82 ± 1,04 \\
\hline
\multirow{7}{*}{\textbf{Texas}} & Community & 34,22 ± 0,43 & 35,99 ± 0,93 & \textbf{36,97 ± 0,81} \\
 & Entertainment & 25,56 ± 1,56 & \textbf{28,30 ± 1,37} & 27,59 ± 1,64 \\
 & Food & 29,56 ± 4,10 & 47,48 ± 2,42 & \textbf{48,67 ± 0,66} \\
 & Nightlife & 25,46 ± 1,13 & 26,06 ± 1,36 & \textbf{26,14 ± 2,66} \\
 & Outdoors & 23,02 ± 1,03 & \textbf{24,62 ± 0,46} & 23,04 ± 1,44 \\
 & Shopping & 38,79 ± 4,12 & \textbf{45,40 ± 1,30} & 43,18 ± 2,19 \\
 & Travel & 29,71 ± 1,32 & \textbf{34,26 ± 0,90} & 33,05 ± 1,51 \\
\hline
\end{tabular}}
\end{table}

\subsection{Predição do Próximo POI}

Os resultados da tarefa de predição do próximo POI são apresentados na Tabela \ref{table:next}. Diferentemente do que ocorre na classificação, essa tarefa apresenta um cenário mais heterogêneo, mas ainda favorável às variantes propostas: os modelos espaço-temporais superam o MTLNet original em 16 das 21 combinações avaliadas. Os ganhos mais expressivos ocorrem em \textit{Food}, categoria em que ambas as variantes superam o \textit{baseline} nos três estados. O melhor resultado é obtido pelo Sphere2Vec-M no Texas, com 48,67 $\pm$ 0,66, enquanto na Califórnia o mesmo modelo alcança 51,28 $\pm$ 0,57, superando com ampla margem o DGI. Melhorias consistentes também aparecem em \textit{Shopping} e \textit{Community}, sugerindo que a combinação entre informação espacial contínua e contexto temporal favorece a modelagem de transições recorrentes entre categorias.

O padrão de co-localização mostrado na Figura \ref{fig:distribuicao_espacial} ajuda a interpretar esse comportamento, pois a presença de agrupamentos densos de categorias como \textit{Food} e \textit{Shopping} favorece transições locais frequentes ao longo da trajetória do usuário. Ainda assim, o \textit{baseline} mantém vantagem em alguns casos, especialmente em \textit{Travel} na Flórida e na Califórnia, além de categorias específicas na Califórnia, como \textit{Entertainment}, \textit{Nightlife} e \textit{Outdoors}. Esses resultados sugerem que deslocamentos de longa distância e relações topológicas entre regiões ainda são melhor representados pelo DGI, cujo grafo espacial captura conexões estruturais que os \textit{encoders} baseados em coordenadas não modelam diretamente. Mesmo com essa limitação, o desempenho geral indica que as representações propostas são mais eficazes na maioria dos cenários de predição sequencial.
